\section{Introduction}
\label{sec:intro}

In the last three years, pretrained forecasting models have moved from research curiosities to
default components. TimesFM \citep{das2024timesfm}, the Chronos family
\citep{ansari2024chronos,ansari2025chronos2}, Sundial \citep{liu2025sundial}, Timer
\citep{liu2024timer}, Moirai \citep{woo2024moirai}, MOMENT \citep{goswami2024moment}, Lag-Llama
\citep{rasul2023lagllama}, Time-MoE \citep{shi2025timemoe} and Tiny Time Mixers
\citep{ekambaram2024ttm} are all released as open checkpoints and are used zero-shot on series they
never saw. Their architectures differ a great deal: decoder-only stacks over 32-step patches,
encoder--decoders over quantized single steps, encoders with cross-series attention, flow-matching
heads that sample. Leaderboards say which is more accurate on which benchmark. They do not say
\emph{what these models represent}, whether architectures this different end up representing the same
things, or whether any internal structure we can name actually drives the forecast.

Mechanistic interpretability has produced a mature toolkit for exactly these questions in language
models: probes \citep{alain2017probes,hewitt2019control}, representational similarity
\citep{kornblith2019cka}, activation patching \citep{vig2020causal,meng2022rome,zhang2024patching},
lenses \citep{nostalgebraist2020logitlens,belrose2023tuned} and sparse dictionaries
\citep{bricken2023monosemanticity,cunningham2024sae,gao2025topk}, packaged in libraries such as
TransformerLens \citep{nanda2022transformerlens}, nnsight \citep{fiottokaufman2025nnsight}, pyvene
\citep{wu2024pyvene} and SAELens \citep{bloom2024saelens}. Early work has begun to apply pieces of it
to TSFMs \citep{wilinski2025exploring,mishra2026dissecting,oublal2026timesae,divo2026zeitgeist}.
Three properties of forecasting models, however, break the direct transfer of these tools, and each of
them causes silent errors rather than crashes:

\begin{enumerate}[label=(\roman*)]
  \item \textbf{Tokens are not comparable across models.} One token is one time step in Chronos-T5,
    32 steps in TimesFM, 16 in Chronos-2 and 96 in Timer. ``Layer 6, position 10'' refers to
    different stretches of time in different models, so every cross-model comparison needs an explicit,
    verified mapping from tokens to time.
  \item \textbf{There is no ground truth for concepts.} Text comes with linguistic annotations. Real
    series do not come with labels for their trend, seasonal period, changepoints or noise process,
    so a probe that ``finds seasonality'' cannot be checked against the answer.
  \item \textbf{Effects are small and easy to fake.} Removing a single sparse feature typically moves a
    forecast by a few hundredths of the context's standard deviation, less than the reconstruction
    error of the dictionary itself (Section~\ref{sec:rq3}). At that scale, a precision mismatch, a
    patch the head never reads, or a null measured on the wrong rows produces a clean, confident and
    wrong result.
\end{enumerate}

This paper presents \tool{}, a system designed around these three problems, and uses it to study
three research questions:

\begin{description}[leftmargin=1.2em,font=\normalfont\bfseries]
  \item[RQ1 (concepts).] What interpretable, forecast-relevant concepts do TSFMs learn?
  \item[RQ2 (sharing).] Are these concepts shared across models with different architectures, and in
    what sense: do the models respond to the same inputs, have the same effect, or both?
  \item[RQ3 (causality).] Do the concepts causally affect the forecast, and how much of the forecast do
    they account for?
\end{description}
A supporting question, \textbf{RQ4}, asks where in depth these computations happen and how far such
depth claims can be compared across architectures.

\paragraph{Approach.} \tool{} rests on three commitments that together distinguish it from existing
tooling (Table~\ref{tab:tools}):
\emph{controlled} measurement, in which a synthetic benchmark with exact generative ground truth and a
battery of structured perturbations gives every probe a known answer to be validated against;
\emph{causal} concepts, in which a feature counts only if removing it changes the forecast by more
than removing a random direction of the same size does, on the same series, with reach controls
proving the intervention landed; and \emph{confirmatory} testing, in which dev-corpus findings are
frozen in a hash-pinned registry and tested once on a sealed private split. Two further design choices
make these commitments usable: architecture-agnostic \emph{time alignment}, which measures each
model's token-to-time map and pools all models onto common windows, and a \emph{one-command report}
that explains every figure in plain language, states what each number can and cannot support, and
renders every skipped analysis with its reason.

\paragraph{Contributions.}
\begin{enumerate}[label=\textbf{C\arabic*.},leftmargin=2.2em]
  \item \textbf{An open-source system} (\tool{}, ${\sim}$61k lines of analysis code, 171 test modules)
    that goes from a configuration listing model checkpoints to a self-contained HTML report through
    19 fingerprinted stages, with a zero-code adapter path for new Hugging Face checkpoints
    (Section~\ref{sec:architecture}).
  \item \textbf{A controlled benchmark generator} with two explicit leakage tiers, a banded-DTW leakage
    audit, sealed dev/private splits with fresh epochs, and bit-exact regeneration
    (Section~\ref{sec:benchmark}).
  \item \textbf{A causal concept pipeline} that defines concepts by their effect on nine forecast
    channels, pools them across models into an atlas, and separates two notions of sharing that prior
    work conflates: \emph{same inputs} (top-$k$ overlap against a stratum-matched null) and
    \emph{same effect on the same inputs} (joint ablation against activation-matched floors, with a
    lower-tail test for disagreement) (Section~\ref{sec:concepts}).
  \item \textbf{A statistical protocol} in which the series is always the resampling unit, every
    comparison enters a multiplicity ledger whose attainable minimum $p$-value is checked, and
    confirmation happens once, on sealed data (Section~\ref{sec:statistics}).
  \item \textbf{An empirical study} of four TSFMs that answers RQ1--RQ4 with typed evidence and a
    held-out replication (Section~\ref{sec:casestudy}).
  \item \textbf{A catalogue of measurement failures}: fourteen documented cases in which an
    uncontrolled analysis produced a plausible, well-formed and wrong result, each paired with the
    built-in control that caught it (Section~\ref{sec:controls}). We argue this is the most transferable
    contribution: it shows concretely why interpretability pipelines for forecasting models need these
    controls.
\end{enumerate}

\paragraph{Main findings.} On TimesFM~2.5, Chronos-2, Sundial and Chronos-Bolt:
\begin{itemize}
  \item \textbf{RQ1.} Causal sparse features are real (293 features across 29 layers clear a row-matched
    random-direction null) but overwhelmingly act on forecast \emph{level} and \emph{shape}. They tile a
    low-dimensional effect space continuously rather than forming discrete concepts: nine readable
    ``families'' cover 98\% of them but do not beat a column-shuffle null.
  \item \textbf{RQ2.} The models agree on \emph{which series} a concept groups: all 20 registered
    cross-model transfer claims replicate on 800 sealed series (top-$k$ AUC 0.971--0.995). They rarely
    agree on \emph{what the concept does}: only 11 of 1{,}223 shared-input tests show the same causal
    effect, the dominant sharing class is ``convergent'' (same effect, different inputs), and no concept
    has causal members in all four models. Seven pairwise similarity metrics rank the six model pairs
    inconsistently (Kendall's $W=0.28$), and representational similarity cannot detect lineage.
  \item \textbf{RQ3.} Single features are causally load-bearing but small: the median feature effect is
    about 3.5 times smaller than the effect of replacing a layer by its SAE reconstruction, and 64\% of
    visually prominent features have no measurable causal effect at all.
  \item \textbf{RQ4.} Where in depth models react to structured corruptions replicates on held-out data
    (9 of 9 registered corruptions for the TimesFM/Chronos-2 pair), but cross-architecture depth claims
    are qualified by how much of each model's forward pass is actually observed (from 14\% to 98\% of
    forecast FLOPs).
\end{itemize}

\paragraph{How to read this paper.} Sections~\ref{sec:principles}--\ref{sec:statistics} describe the
system; Section~\ref{sec:casestudy} uses it; Section~\ref{sec:controls} shows what its controls
prevent. We do not attempt to report every result the system has produced. The full ledger of
findings, each with exact numbers, evidence class, confirmation status and a reproduction pointer,
is \repo{FINDINGS.md}; Appendix~\ref{app:ledger} indexes the entries cited here. The complete example
run analyzed in Section~\ref{sec:casestudy}, including its report and every per-stage artifact, is
\repodir{examples/concept_atlas_v2}, and every data figure in this paper is regenerated from it by
\repo{paper/figures/make_figures.py}.

\begin{figure}[t]
\centering
\begin{tikzpicture}[
  node distance=4mm and 6mm,
  box/.style={draw, rounded corners=2pt, align=center, font=\footnotesize, minimum height=8mm, fill=white},
  pkg/.style={draw=gray!60, dashed, rounded corners=4pt, inner sep=6pt},
  lab/.style={font=\scriptsize\itshape, text=gray!70!black},
  arr/.style={-{Stealth[length=2mm]}, thick, gray!70!black}]
  % benchmark package
  \node[box] (gen) {Generators\\{\scriptsize parametric, archetypes,}\\{\scriptsize real-derived}};
  \node[box, below=of gen] (aud) {Leakage audit\\{\scriptsize banded DTW vs.\ references}};
  \node[box, below=of aud] (seal) {Seal \& split\\{\scriptsize sha256 digests, epochs}};
  \node[box, below=of seal] (val) {Validation\\{\scriptsize catch22 diversity, exchangeability}};
  \draw[arr] (gen) -- (aud); \draw[arr] (aud) -- (seal); \draw[arr] (seal) -- (val);
  \begin{scope}[on background layer]
    \node[pkg, fit=(gen)(val), label={[font=\small\bfseries]above:\texttt{tsfm\_benchmark}}] (bp) {};
  \end{scope}
  % corpora
  \node[box, right=12mm of aud, fill=soft] (dev) {Dev split\\{\scriptsize ground-truth labels}};
  \node[box, below=of dev, fill=soft] (priv) {Private split\\{\scriptsize sealed, used once}};
  \draw[arr] (seal.east) -- ++(4mm,0) |- (dev.west);
  \draw[arr] (seal.east) -- ++(4mm,0) |- (priv.west);
  % lens package
  \node[box, right=14mm of dev, yshift=9mm] (ad) {Adapters \& tiers\\{\scriptsize measured spans}};
  \node[box, below=of ad] (st) {19 fingerprinted stages\\{\scriptsize behavior $\to$ geometry $\to$ causality $\to$ concepts}};
  \node[box, below=of st] (reg) {Register $\to$ confirm once};
  \node[box, below=of reg] (rep) {Self-contained HTML report\\{\scriptsize typed findings, plain-language notes}};
  \draw[arr] (ad) -- (st); \draw[arr] (st) -- (reg); \draw[arr] (reg) -- (rep);
  \begin{scope}[on background layer]
    \node[pkg, fit=(ad)(rep), label={[font=\small\bfseries]above:\texttt{tsfm\_lens}}] (lp) {};
  \end{scope}
  \draw[arr] (dev.east) -- (st.west);
  \draw[arr] (priv.east) -- (reg.west);
  \node[box, above=5mm of ad, fill=soft] (ck) {Model checkpoints + one YAML config};
  \draw[arr] (ck) -- (ad);
\end{tikzpicture}
\caption{System overview. \texttt{tsfm\_benchmark} produces leakage-audited, sealed corpora with
exact generative ground truth; \texttt{tsfm\_lens} takes a configuration naming model checkpoints
and produces one report. Dev-split findings are exploratory; the private split is opened once, by
the \texttt{confirm} stage, for claims registered in advance. The two packages install separately
because their dependencies barely overlap.}
\label{fig:overview}
\end{figure}
